\section{Related Work}

\paragraph{Instruction-data selection.}
Instruction-data selectors usually estimate quality or diversity independently of the target PEFT configuration. LIMA supports instruction tuning with 1,000 curated examples~\cite{zhou2023lima}. Model-based filters instead use judgments or instruction properties: AlpaGasus filters Alpaca examples with LLM judgments~\cite{chen2024alpagasus}, InsTag measures semantic diversity and complexity through instruction tags~\cite{lu2023instag}, and DEITA combines quality, complexity, and diversity~\cite{liu2024deita}. Other methods select for explicit diversity through quality--diversity trade-offs, log-determinant objectives, or partition-aware sampling~\cite{bukharin2023diversity, wang2024diversity, rao2024commonit}. These methods improve data efficiency, but their rankings omit the target PEFT's trainable parameters.

\paragraph{Proxy-model and influence selection.}
Proxy-model selectors provide reusable rankings without computing gradients for every target PEFT. Cherry LLM uses instruction-following difficulty~\cite{li2023cherry}; Superfiltering studies when weak models can rank data for stronger models~\cite{wang2024superfiltering}; and SmallToLarge transfers training-trajectory information from small models to larger targets~\cite{yang2024smalltolarge}. Because these scores encode only example and task information, a PEFT change leaves the ranking unchanged.

Influence-based selectors model training effects more directly, but their gradients tie the score to a target parameterization. Classical influence functions use second-order approximations~\cite{koh2017understanding}, while TracIn accumulates gradient similarities across checkpoints~\cite{pruthi2020tracin}. LESS builds a low-dimensional gradient datastore for instruction selection~\cite{xia2024less}, and DataInf specializes influence approximation for LoRA-tuned generative models~\cite{kwon2023datainf}. Preserving configuration specificity therefore requires rebuilding PEFT-dependent gradients when the target changes. PCU-Select instead conditions one scorer on a structured PEFT code, combining reusable sample signals with a target-specific ranking. Our evaluation compares PCU-Select with both a per-PEFT LESS implementation and an Influence baseline that uses one shared datastore.

\paragraph{Parameter-efficient fine-tuning.}
PEFT methods reduce adaptation cost in different ways, so they expose different trainable parameters. Adapters insert bottleneck modules into Transformer layers~\cite{houlsby2019adapters, pfeiffer2020adapterfusion}, while prefix and prompt tuning optimize continuous prefix states or prompts~\cite{li2021prefixtuning, lester2021prompttuning}. LoRA learns low-rank additive updates~\cite{hu2022lora}, and QLoRA combines quantization with LoRA to reduce memory use~\cite{dettmers2023qlora}. BitFit tunes bias terms~\cite{zaken2022bitfit}, whereas (IA)$^3$ rescales attention and feed-forward activations~\cite{liu2022ia3}. Prior PEFT studies usually compare accuracy, parameter count, memory, or inference overhead while holding the training data fixed. We study the complementary question: whether a fixed candidate pool has the same example ranking across PEFT configurations.

\paragraph{Joint data- and parameter-efficient adaptation.}
The closest methods use PEFT gradients for data selection, but their utility estimates remain specific to one target configuration. DataInf and LESS compute selection signals from LoRA-style training~\cite{kwon2023datainf, xia2024less}. Geometry-aware targeted selection further shows that useful data depends on the directions available to PEFT optimization~\cite{min2026gist}. These methods must regenerate configuration-dependent signals when the target changes.

PCU-Select instead describes each configuration by its intervention sites, update type, capacity, and training recipe. This representation lets one scorer reuse sample and task signals while still producing a different ranking for each configuration.
